% !TEX root = ./signals_and_systems_notes_en.tex
% !TEX program = xelatex
% English edition of 信号与系统笔记.tex — all display math kept unchanged.
% Layout (fonts / cover / boxes / TOC style) lives in _preamble_en.tex
% !TEX root = ./signals_and_systems_notes_en.tex
% !TEX program = xelatex
%==================================================================
%  Signals and Systems Notes — layout template (English edition)
%  Compile: run xelatex twice (or latexmk -xelatex)
%  Based on _preamble.tex; only the language-specific strings
%  (chapter label, box names, cover, headers) were changed.
%==================================================================
\documentclass[UTF8,11pt,oneside,openany,scheme=chinese,fontset=none]{ctexbook}

% ---------- page ----------
\usepackage[a4paper,left=2.8cm,right=2.8cm,top=2.7cm,bottom=2.7cm,
            headheight=15pt,headsep=11pt,footskip=16pt]{geometry}

% ---------- math & graphics ----------
\usepackage{amsmath,amssymb,amsthm}
\usepackage{newtxmath}                    % formulas: Times-flavoured
\usepackage{tikz}
\usetikzlibrary{calc,positioning}
\usepackage{pgfplots}
\pgfplotsset{compat=1.18}
\usepackage{forest}
\usepackage{graphicx}
\usepackage{array,makecell,booktabs,tabularx,multirow}
\usepackage{enumitem}
\usepackage[most]{tcolorbox}
\usepackage{etoolbox}
\usepackage{fancyhdr}
\usepackage{titlesec}
\usepackage{titletoc}
\usepackage[font=small,labelfont={sf,bf},labelsep=quad,
            justification=centering,skip=6pt]{caption}
\usepackage{subcaption}                   % side-by-side subfigures
\usepackage{float}                        % for [H]
\usepackage[hidelinks,bookmarksnumbered]{hyperref}

% ---------- colors ----------
\definecolor{accent}{HTML}{2E4A63}   % main: deep indigo
\definecolor{warm}{HTML}{B4562C}     % secondary: terracotta
\definecolor{ink}{HTML}{1F2A33}      % body black
\definecolor{muted}{HTML}{6B7B8A}    % muted grey-blue
\definecolor{rule}{HTML}{D5DDE5}     % thin rules
\definecolor{cool}{HTML}{F1F4F8}     % definition box
\definecolor{cream}{HTML}{FBF5EC}    % note box
\definecolor{soft}{HTML}{F6F2EA}     % code background
\definecolor{warnbg}{HTML}{FBEDE6}   % warning box

% ---------- fonts ----------
\setmainfont{TeX Gyre Termes}
\setsansfont{TeX Gyre Heros}[Scale=0.92]
\setmonofont{Consolas}[Scale=0.85]
% CJK fonts retained (harmless if no Chinese appears in the body)
\setCJKmainfont{LXGW Neo ZhiSong}[AutoFakeBold=2.2,AutoFakeSlant=0.15]
\setCJKsansfont{HarmonyOS Sans SC}[AutoFakeBold=2.2,AutoFakeSlant=0.15]
\setCJKmonofont{HarmonyOS Sans SC}[AutoFakeBold=2.0]

% ---------- notation ----------
\newcommand{\dd}{\mathrm{d}}
\newcommand{\lrarr}{\;\longleftrightarrow\;}
\DeclareMathOperator{\Sa}{Sa}
\DeclareMathOperator{\Si}{Si}
\DeclareMathOperator{\sgn}{sgn}
\newcommand{\highlight}[1]{\colorbox{yellow!28}{\textbf{#1}}}
\allowdisplaybreaks[2]

% ---------- English names for ctex-managed strings ----------
\renewcommand{\proofname}{Proof}
\renewcommand{\figurename}{Figure}
\renewcommand{\tablename}{Table}
\renewcommand{\contentsname}{Contents}
\ctexset{
  chapter/number=\arabic{chapter},   % TOC: arabic number
  chapter/name={,},                  % remove "第 … 章" wrapper
}

% ---------- body tweaks ----------
\setlength{\parskip}{2pt plus 1pt}
\renewcommand{\arraystretch}{1.25}
\setlength{\abovedisplayskip}{7pt plus 2pt minus 2pt}
\setlength{\belowdisplayskip}{7pt plus 2pt minus 2pt}
\setlength{\abovedisplayshortskip}{4pt}
\setlength{\belowdisplayshortskip}{5pt}
\clubpenalty=1000 \widowpenalty=1000 \displaywidowpenalty=1000

%==================================================================
%  Sectioning
%==================================================================
% chapter: vertical bar + small caps label + big title + thin rule
\titleformat{\chapter}[display]
  {\normalfont\raggedright}
  {\vspace*{-10pt}{\color{accent}\rule[-2.5pt]{3pt}{13pt}}\hspace{7pt}%
   {\sffamily\bfseries\color{accent}\fontsize{11}{11}\selectfont CHAPTER\,\thechapter}}
  {0.6em}
  {\sffamily\bfseries\color{ink}\fontsize{25}{31}\selectfont}
  [\vspace{0.7em}{\color{rule}\hrule height 0.7pt}]
\titlespacing*{\chapter}{0pt}{10pt}{26pt}

% unnumbered chapter (TOC, preface)
\titleformat{name=\chapter,numberless}[display]
  {\normalfont\raggedright}
  {\vspace*{-10pt}}
  {0pt}
  {\sffamily\bfseries\color{ink}\fontsize{24}{30}\selectfont}
  [\vspace{0.7em}{\color{rule}\hrule height 0.7pt}]
\titlespacing*{name=\chapter,numberless}{0pt}{10pt}{26pt}

% section
\titleformat{\section}[block]
  {\normalfont\Large\sffamily\bfseries\color{ink}}
  {\raisebox{1.2pt}{\colorbox{accent}{\strut\hspace{2pt}\color{white}\thesection\hspace{2pt}}}}{0.7em}{}
  [\vspace{-0.15em}{\color{rule}\hrule height 0.7pt}\vspace{-0.1em}]
\titlespacing*{\section}{0pt}{1.5em}{0.9em}

% subsection
\titleformat{\subsection}[block]
  {\normalfont\large\sffamily\bfseries\color{ink}}
  {\color{warm}\thesubsection}{0.6em}{}
\titlespacing*{\subsection}{0pt}{1.3em}{0.6em}

% subsubsection
\titleformat{\subsubsection}[block]
  {\normalfont\normalsize\sffamily\bfseries\color{ink}}
  {\color{warm}$\blacktriangleright$}{0.45em}{}
\titlespacing*{\subsubsection}{0pt}{1.1em}{0.5em}

%==================================================================
%  Headers / footers
%==================================================================
\renewcommand{\chaptermark}[1]{\markboth{\thechapter\enspace #1}{}}
\renewcommand{\sectionmark}[1]{\markright{\thesection\enspace #1}}

\pagestyle{fancy}
\fancyhf{}
\fancyhead[L]{\sffamily\footnotesize\color{muted}Signals and Systems Notes}
\fancyhead[R]{\sffamily\footnotesize\color{muted}\nouppercase{\leftmark}}
\fancyfoot[C]{\sffamily\footnotesize\color{muted}\thepage}
\renewcommand{\headrulewidth}{0.4pt}
\renewcommand{\headrule}{\hbox to\headwidth{\color{rule}\leaders\hrule height \headrulewidth\hfill}}
\fancypagestyle{plain}{%
  \fancyhf{}%
  \fancyfoot[C]{\sffamily\footnotesize\color{muted}\thepage}%
  \renewcommand{\headrulewidth}{0pt}}

%==================================================================
%  Table of contents
%==================================================================
\renewcommand{\contentsname}{Contents}
\titlecontents{chapter}[0em]
  {\addvspace{7.5pt}\sffamily\bfseries\color{ink}}
  {\contentslabel{5.6em}\color{accent}}
  {\hspace*{-5.6em}}
  {\hfill\hspace*{0.45em}\makebox[2.4em][r]{\color{accent}\thecontentspage}}
\titlecontents{section}[3.0em]
  {\addvspace{1.2pt}\small\color{ink}}
  {\contentslabel{2.4em}\color{muted}}
  {\hspace*{-2.4em}}
  {\titlerule*[0.6pc]{.}\hspace*{0.45em}\makebox[2.4em][r]{\color{muted}\thecontentspage}}
\titlecontents{subsection}[3.4em]
  {\addvspace{0.4pt}\small\color{muted}}
  {\contentslabel{3.2em}}
  {\hspace*{-3.2em}}
  {\titlerule*[0.6pc]{.}\hspace*{0.45em}\makebox[2.4em][r]{\color{muted}\thecontentspage}}

%==================================================================
%  Colored boxes (Definition / Note / Warning / Proof)
%==================================================================
\tcbset{
  notebox/.style={
    enhanced, breakable, boxrule=0pt, arc=2pt,
    left=11pt, right=11pt, top=8pt, bottom=8pt,
    fontupper=\normalsize, before skip=9pt, after skip=9pt,
    segmentation style={solid, rule, line width=0.5pt}
  }
}

% —— Definition (numbered) ——
\newcounter{definition}[chapter]
\newenvironment{definition}[1][]{%
  \refstepcounter{definition}\par\addvspace{2pt}%
  \begin{tcolorbox}[notebox,colback=cool,borderline west={2.4pt}{0pt}{accent}]%
  \noindent{\sffamily\bfseries\color{accent}Definition \thechapter.\arabic{definition}}%
  \ifstrempty{#1}{}{\ \ {\sffamily\color{accent}(#1)}}%
  \par\nobreak\vspace{3pt}\nobreak\ignorespaces
}{%
  \end{tcolorbox}%
}

% —— Definition (unnumbered) ——
\newenvironment{definition*}[1][]{%
  \par\addvspace{2pt}%
  \begin{tcolorbox}[notebox,colback=cool,borderline west={2.4pt}{0pt}{accent}]%
  \noindent{\sffamily\bfseries\color{accent}Definition}%
  \ifstrempty{#1}{}{\ \ {\sffamily\color{accent}(#1)}}%
  \par\nobreak\vspace{3pt}\nobreak\ignorespaces
}{%
  \end{tcolorbox}%
}

% —— Note / example (cream) ——
\newenvironment{note}[1][]{%
  \par\addvspace{2pt}%
  \begin{tcolorbox}[notebox,colback=cream,borderline west={2.4pt}{0pt}{warm}]%
  \ifstrempty{#1}{}{\noindent{\sffamily\bfseries\color{warm}#1}\par\vspace{3pt}}%
  \ignorespaces
}{%
  \end{tcolorbox}%
}

% —— Warning (terracotta) ——
\newenvironment{warning}[1][]{%
  \par\addvspace{2pt}%
  \begin{tcolorbox}[notebox,colback=warnbg,borderline west={2.4pt}{0pt}{warm}]%
  \noindent{\sffamily\bfseries\color{warm}$\blacktriangleright$\,Warning}%
  \ifstrempty{#1}{}{\ \ {\sffamily\color{warm}--- #1}}%
  \par\nobreak\vspace{3pt}\nobreak\ignorespaces
}{%
  \end{tcolorbox}%
}

% —— Theorem / Example (reserved) ——
\newcounter{theorem}[chapter]
\newenvironment{theorem}[1][]{%
  \refstepcounter{theorem}\par\addvspace{2pt}%
  \begin{tcolorbox}[notebox,colback=white,borderline west={2.4pt}{0pt}{accent},
                    borderline north={0.5pt}{0pt}{rule},borderline south={0.5pt}{0pt}{rule}]%
  \noindent{\sffamily\bfseries\color{accent}Theorem \thechapter.\arabic{theorem}}%
  \ifstrempty{#1}{}{\ \ {\sffamily\color{accent}(#1)}}%
  \par\nobreak\vspace{3pt}\nobreak\ignorespaces
}{\end{tcolorbox}}

\newenvironment{example}[1][]{%
  \par\addvspace{2pt}%
  \begin{tcolorbox}[notebox,colback=white,borderline west={2.4pt}{0pt}{muted}]%
  \noindent{\sffamily\bfseries\color{muted}Example}%
  \ifstrempty{#1}{}{\ \ {\sffamily\color{muted}#1}}%
  \par\nobreak\vspace{3pt}\nobreak\ignorespaces
}{\end{tcolorbox}}

% —— Proof (customized amsthm) ——
\renewenvironment{proof}[1][\proofname]{%
  \par\addvspace{4pt}%
  \pushQED{\qed}%
  \normalfont\topsep2pt\trivlist
  \item[\hspace*{0pt}{\sffamily\bfseries\color{accent}#1}]\hspace*{0.4em}\ignorespaces
}{%
  \popQED\endtrivlist\@endpefalse
}

%==================================================================
%  Code listings (reserved)
%==================================================================
\usepackage{listings}
\lstdefinestyle{plain}{
  backgroundcolor=\color{soft}, basicstyle=\ttfamily\small,
  keywordstyle=\color{accent}\bfseries, commentstyle=\color{muted}, stringstyle=\color{warm},
  frame=none, framesep=0pt, xleftmargin=22pt, xrightmargin=6pt,
  numbers=left, numberstyle=\ttfamily\scriptsize\color{muted}, numbersep=8pt,
  breaklines=true, showstringspaces=false, tabsize=2,
  aboveskip=9pt, belowskip=9pt,
  captionpos=t, extendedchars=false
}
\lstset{style=plain}
\renewcommand{\lstlistingname}{Listing}

%==================================================================
%  Cover
%==================================================================
\makeatletter
\newcommand{\subtitle}[1]{\gdef\@subtitle{#1}}
\gdef\@subtitle{}
\newcommand{\institute}[1]{\gdef\@institute{#1}}
\gdef\@institute{}
\newcommand{\cover}[1]{\gdef\@coverimage{#1}}
\gdef\@coverimage{cover.png}
\newcommand{\extrainfo}[1]{\gdef\@extrainfo{#1}}
\gdef\@extrainfo{}
\renewcommand{\maketitle}{%
  \begin{titlepage}%
    \begin{tikzpicture}[remember picture,overlay]
      \node[anchor=north west,inner sep=0pt] at (current page.north west)
        {\includegraphics[width=\paperwidth]{\@coverimage}};
      \node[anchor=north west,align=left,text width=15.4cm,inner sep=0pt]
        at ([xshift=2.8cm,yshift=-15.2cm]current page.north west) {%
          {\color{accent}\rule[-2pt]{2.6pt}{12pt}}\hspace{7pt}%
          {\sffamily\bfseries\color{accent}\fontsize{10.5}{10.5}\selectfont S\,T\,U\,D\,Y\ \ N\,O\,T\,E\,S}\par
          \vspace{16pt}
          {\sffamily\bfseries\color{ink}\fontsize{34}{40}\selectfont \@title}\par
          \vspace{10pt}
          {\sffamily\color{muted}\fontsize{13}{18}\selectfont \@subtitle}\par
          \vspace{26pt}
          {\color{accent}\rule{4.2cm}{1.6pt}}\par
          \vspace{14pt}
          {\sffamily\color{ink}\fontsize{11}{14}\selectfont \@author}\par
          \vspace{3pt}
          {\sffamily\color{muted}\fontsize{10}{13}\selectfont \@institute}\par
        };
      \node[anchor=south west,align=left,text width=15.4cm,inner sep=0pt]
        at ([xshift=2.8cm,yshift=3.0cm]current page.south west) {%
          {\color{rule}\rule{\linewidth}{0.5pt}}\par\vspace{6pt}
          {\sffamily\color{muted}\fontsize{9}{12}\selectfont \@date\hfill\@extrainfo}\par
        };
    \end{tikzpicture}%
  \end{titlepage}%
  \setcounter{page}{1}%
}
\makeatother

\graphicspath{{figures/}}

% ---------- cover data ----------
\title{Signals and Systems Notes}
\subtitle{Time-Domain and Frequency-Domain Analysis of Continuous and Discrete Systems}
\author{Zhang Junxiang}
\institute{Communication University of China}
\date{October 2026}
\cover{cover.png}
\extrainfo{The starting point of the journey}

\usepackage{subcaption}                   % side-by-side subfigures

\begin{document}
\maketitle
\frontmatter
\tableofcontents
\mainmatter
\clearpage

% ==================================================================
\chapter{Introduction}
% ==================================================================

% \section{Signals and Systems}

% \begin{definition*}
% A signal is a representation of information; information is the content
% carried by a signal.

% A system is an aggregate of interacting, interdependent things that
% together perform a specific function.
% \end{definition*}

% \section{Description, Classification, and Examples of Signals}

% \begin{figure}[htbp]
% \centering
% \resizebox{\linewidth}{!}{%
% \begin{forest}
%     for tree={
%         grow'=0,
%         parent anchor=east,
%         child anchor=west,
%         edge path={
%             \noexpand\path[\forestoption{edge}] (!u.parent anchor) -- +(5pt,0) |- (.child anchor)\forestoption{edge label};
%         },
%         l sep=12pt,
%         s sep=6pt,
%         font=\small,
%         align=center,
%         draw,
%         rounded corners,
%     }
%     [Classification of signals,
%         [deterministic?,
%             [deterministic signals]
%             [random signals]
%         ]
%         [regularity,
%             [{periodic signals\\$f(t)=f(t+nT)$}]
%             [aperiodic signals]
%         ]
%         [discrete vs.\ continuous,
%             [continuous-time signals]
%             [discrete-time signals,
%                 [sampled signals\\(continuous amplitude)]
%                 [{digital signals\\(amplitude 0 or 1)}]
%             ]
%         ]
%         [dimension,
%             [one-dimensional signals]
%             [multi-dimensional signals]
%         ]
%     ]
% \end{forest}}
% \caption{Classification of signals}
% \end{figure}

% Signals can be classified from many viewpoints, and the classes are not
% mutually exclusive: one and the same signal may be deterministic, periodic,
% continuous-time, and one-dimensional all at once.

% \section{Typical Signals}

% \begin{enumerate}
%     \item Exponential signals
%     \[
%         f(t)=ke^{at}
%     \]
%     with the time constant $\tau$ defined as
%     \[
%         \tau=\frac{1}{|a|}.
%     \]

%     \begin{figure}[htbp]
%     \centering
%     \begin{tikzpicture}
%         \begin{axis}[
%             width=10cm, height=6cm,
%             xmin=0, xmax=3.5,
%             ymin=0, ymax=8,
%             xlabel={$t$}, ylabel={$f(t)$},
%             axis lines=center,
%             grid=major,
%             grid style={dashed, gray!50},
%             xtick={0,1,2,3},
%             ytick={0,2,4,6,8},
%             samples=100,
%             legend pos=north west,
%             domain=0:3,
%         ]
%             \addplot[blue, thick] {2^x};
%             \addlegendentry{$a>1$ (increasing)}

%             \addplot[red, thick, dashed] {0.5^x};
%             \addlegendentry{$0<a<1$ (decreasing)}
%         \end{axis}
%     \end{tikzpicture}
%     \caption{Exponential signals}
%     \end{figure}

%     \textcolor{red}{\textbf{The integral and the derivative of an exponential
%     signal are still exponential signals.}}

%     \item Sinusoidal signals
%     \[
%         f(t)=A\sin(\omega t+\theta)
%     \]
%     This is an ideal, undamped sinusoid; with damping one multiplies by an
%     exponentially decaying factor, i.e.\ the damped sinusoid
%     $Ke^{-\alpha t}\sin(\omega t+\theta)$.

%     \item Complex exponential signals
%     \[
%         f(t)=Ke^{st},\qquad s=\sigma+j\omega
%     \]
%     Using Euler's formula,
%     \[
%         f(t)=Ke^{(\sigma+j\omega)t}=Ke^{\sigma t}e^{j\omega t}
%         =Ke^{\sigma t}\bigl[\cos(\omega t)+j\sin(\omega t)\bigr].
%     \]
%     Thus $\sigma$ controls the decay or growth of the envelope and $\omega$
%     the rate of oscillation; taking the real part yields the damped sinusoid,
%     so the complex exponential is the basic building block for describing
%     LTI systems.

%     \item The sampling function
%     \[
%         \Sa(t)=\frac{\sin t}{t}
%     \]
%     Main properties:
%     \begin{itemize}
%         \item $\Sa(t)$ is even, with maximum value $1$ at $t=0$;
%         \item $\Sa(k\pi)=0\ (k=\pm1,\pm2,\dots)$;
%         \item $\displaystyle\int_{0}^{\infty}\Sa(t)\,\dd t=\frac{\pi}{2}$,
%               $\displaystyle\int_{-\infty}^{\infty}\Sa(t)\,\dd t=\pi$.
%     \end{itemize}

%     \begin{proof}
%     Let
%     \[
%         I(a)=\int_{0}^{\infty}e^{-at}\frac{\sin t}{t}\,\dd t \qquad (a>0).
%     \]
%     Differentiating with respect to the parameter $a$ and interchanging
%     differentiation and integration:
%     \[
%         I'(a)=-\int_{0}^{\infty}e^{-at}\sin t\,\dd t=-\frac{1}{1+a^{2}}.
%     \]
%     Since $I(\infty)=0$, integrating both sides gives
%     \[
%         I(a)=\int_{a}^{\infty}\frac{\dd u}{1+u^{2}}=\frac{\pi}{2}-\arctan a .
%     \]
%     Letting $a\to 0^{+}$,
%     \[
%         \int_{0}^{\infty}\Sa(t)\,\dd t=\lim_{a\to0^{+}}I(a)=\frac{\pi}{2}.
%     \]
%     Since $\Sa(t)$ is even,
%     $\displaystyle\int_{-\infty}^{\infty}\Sa(t)\,\dd t=\pi$.
%     \end{proof}
% \end{enumerate}

% ==================================================================
\chapter{Time-Domain Analysis of Continuous-Time Systems}
% ==================================================================

\section{The Classical Solution of Differential Equations}

The mathematical model relating the input and the response of an LTI
continuous-time system is an $n$-th order linear differential equation with
constant coefficients,
\[
    \sum_{j=0}^{n}a_j\,y^{(j)}(t)=\sum_{i=0}^{m}b_i\,f^{(i)}(t).
\]
The complete solution of the differential equation consists of a homogeneous
solution and a particular solution,
\[
    y(t)=y_h(t)+y_p(t).
\]
The homogeneous solution solves the homogeneous equation; its form is
determined by the characteristic roots:

\begin{table}[htbp]
\centering
\caption{Forms of the homogeneous solution $y_h(t)$}
\small
\begin{tabular}{@{}p{0.31\linewidth}p{0.61\linewidth}@{}}
\toprule
Characteristic root $\lambda$ & Homogeneous solution $y_h(t)$ \\
\midrule
simple real root & $e^{\lambda t}$ \\
real root of multiplicity $r$ & $(C_{r-1}t^{r-1}+C_{r-2}t^{r-2}+\cdots+C_1t+C_0)\,e^{\lambda t}$ \\
a pair of complex conjugate roots $\lambda_{1,2}=\alpha\pm j\beta$ &
$e^{\alpha t}\bigl(C\cos\beta t+D\sin\beta t\bigr)$ \\
complex conjugate roots of multiplicity $r$ &
$e^{\alpha t}\bigl[A_{r-1}t^{r-1}\cos(\beta t+\theta_{r-1})+\cdots+A_0\cos(\beta t+\theta_0)\bigr]$ \\
\bottomrule
\end{tabular}
\end{table}

The functional form of the particular solution depends on the form of the
input:

\begin{table}[htbp]
\centering
\caption{Forms of the particular solution $y_p(t)$}
\small
\begin{tabular}{@{}p{0.26\linewidth}p{0.40\linewidth}p{0.26\linewidth}@{}}
\toprule
Input $f(t)$ & Particular solution $y_p(t)$ & Condition \\
\midrule
$t^{m}$ & $t^{r}\bigl(P_mt^{m}+P_{m-1}t^{m-1}+\cdots+P_0\bigr)$ &
characteristic root $0$ of multiplicity $r$ \\
$e^{\alpha t}$ & $\bigl(P_rt^{r}+P_{r-1}t^{r-1}+\cdots+P_1t+P_0\bigr)e^{\alpha t}$ &
$\alpha$ is a characteristic root of multiplicity $r$ \\
$\cos\beta t$ or $\sin\beta t$ & $P\cos\beta t+Q\sin\beta t$ &
no characteristic root equals $j\beta$ \\
\bottomrule
\end{tabular}
\end{table}

\begin{note}
``Multiplicity $r$'' in the table is the generic notation: when the input does
not coincide with a characteristic root, take $r=0$; the particular solution
then carries no factor of $t$ and is the ordinary undetermined-coefficient
form.
\end{note}

\section{On the Values at $0_{-}$ and $0_{+}$}

When solving a differential equation by the classical method, the input
$f(t)$ is usually applied to the system at $t=0$, and the solution holds for
$t>0$. The set of initial values needed to fix the solution refers to the
instant $t=0_{+}$. At $t=0_{-}$ the input has not yet been applied, so the
values of the system depend only on the initial state, not on the input. For
some inputs, the response and its derivatives jump when the input is switched
on, so the jump values must be found.

\subsection{The Integration Method}

This method is based on the integral properties of $\delta(t)$,
$\varepsilon(t)$, and $t\varepsilon(t)$, and on their behaviour at the origin:
$\delta(t)$ has an impulse at the origin; its integral $\varepsilon(t)$ jumps
at the origin; and the integral of $\varepsilon(t)$, namely $t\varepsilon(t)$,
is continuous there.

\begin{note}
\textbf{Example}\quad An LTI system is described by
\[
    y''(t)+3y'(t)+2y(t)=2f'(t)+6f(t),
\]
with $y(0_{-})=2$, $y'(0_{-})=0$, and $f(t)=\varepsilon(t)$. Find
$y(0_{+})$ and $y'(0_{+})$.

\textbf{Solution}\quad Only the initial values are computed here.
Substituting the input, $f'(t)=\delta(t)$, gives
\[
    y''(t)+3y'(t)+2y(t)=2\delta(t)+6\varepsilon(t).
\]
Integrating both sides from $0_{-}$ to $0_{+}$:
\[
    \int_{0_{-}}^{0_{+}}y''(t)\,\dd t
    +3\int_{0_{-}}^{0_{+}}y'(t)\,\dd t
    +2\int_{0_{-}}^{0_{+}}y(t)\,\dd t
    =2\int_{0_{-}}^{0_{+}}\delta(t)\,\dd t
    +6\int_{0_{-}}^{0_{+}}\varepsilon(t)\,\dd t .
\]
Inspecting the structure of the equation, $y''(t)$ is the only term that can
contain $\delta(t)$, while $y'(t)$ is bounded near the origin and $y(t)$ is
continuous; hence the last two integrals on the left vanish, and so does the
second term on the right:
\[
    \bigl[y'(0_{+})-y'(0_{-})\bigr]+3\cdot 0+2\cdot 0=2+0 .
\]
Therefore
\[
    y'(0_{+})=2,\qquad y(0_{+})=y(0_{-})=2 .
\]
\end{note}

\subsection{The Undetermined-Coefficient Method}

When the input contains higher-order derivatives of the impulse function, the
integration method cannot reduce them to continuous functions and thus yields
no unique solution, only a system of equations. In that case use the method of
undetermined coefficients.

\begin{note}
\textbf{Example}\quad Let
\[
    y''(t)+2y'(t)+y(t)=\delta''(t)+2\delta(t).
\]
For the equality to hold, the coefficients of each term on both sides must
match. Since $y''(t)$ must contain a $\delta''(t)$ term, set
\[
    y''(t)=a\,\delta''(t)+b\,\delta'(t)+c\,\delta(t)+r_0(t),
\]
and integrate term by term:
\[
    y'(t)=a\,\delta'(t)+b\,\delta(t)+r_1(t),
    \qquad
    y(t)=a\,\delta(t)+r_2(t).
\]
Substituting back into the original equation and matching coefficients:
\[
    \begin{cases}
        a=1,\\
        b+2a=0,\\
        c+2b+a=2,
    \end{cases}
    \quad\Longrightarrow\quad
    a=1,\ b=-2,\ c=3 .
\]
Integrating the impulsive parts of $y'(t)$ and $y''(t)$ then gives the initial
values.
\end{note}

\begin{warning}
Here $r(t)$ denotes the part whose integral over a neighbourhood of the origin
vanishes; it does not affect the computation of the initial values, so only
the coefficients of the singular terms need to be matched.
\end{warning}

\section{Zero-Input Response}

The response caused by the initial state $\{x(0)\}$ alone is called the
zero-input response, denoted $y_{zi}(t)$. With zero input, the right-hand side
of the differential equation vanishes and it becomes homogeneous:
\[
    \sum_{j=0}^{n}a_j\,y_{zi}^{(j)}(t)=0 .
\]
If all characteristic roots are simple, then
\[
    y_{zi}(t)=\sum_{j=1}^{n}C_{zij}\,e^{\lambda_j t}.
\]
Because the input is zero, the initial values are just the system's initial
values, and no jump occurs.

\section{Zero-State Response}

Zero state means the initial values of the system are zero; the response is
caused by the input $f(t)$ alone and is denoted $y_{zs}(t)$:
\[
    \sum_{j=0}^{n}a_j\,y_{zs}^{(j)}(t)=\sum_{i=0}^{m}b_i\,f^{(i)}(t).
\]
If the characteristic roots of the differential equation are all simple, the
zero-state response is
\[
    y_{zs}(t)=\sum_{j=1}^{n}C_{zsj}\,e^{\lambda_j t}+y_p(t),
\]
where $y_p(t)$ is a particular solution of the equation.

\section{The Complete Response}

\begin{equation}
    y(t)=y_{zi}(t)+y_{zs}(t),
\end{equation}
and each derivative obeys the same linear relation
\[
    y^{(n)}(t)=y_{zi}^{(n)}(t)+y_{zs}^{(n)}(t).
\]
In summary, the complete response of an LTI system can be decomposed either
into the zero-input and zero-state responses, or into the natural and forced
responses:
\[
    y(t)=\sum_{j=1}^{n}C_j e^{\lambda_j t}+y_p(t)
        =\sum_{j=1}^{n}C_{zij}e^{\lambda_j t}+\sum_{j=1}^{n}C_{zsj}e^{\lambda_j t}+y_p(t).
\]

\begin{figure}[htbp]
\centering
\resizebox{\linewidth}{!}{%
\begin{forest}
    for tree={
        grow'=0,
        parent anchor=east,
        child anchor=west,
        edge path={
            \noexpand\path[\forestoption{edge}] (!u.parent anchor) -- +(5pt,0) |- (.child anchor)\forestoption{edge label};
        },
        l sep=10pt,
        s sep=4pt,
        inner sep=1.5pt,
        font=\small,
        align=center,
        draw,
        rounded corners,
    }
    [Complete response $y(t)$
        [Zero-input response $y_{zi}(t)$
            [determined by the\\initial state alone]
            [{input is zero, no jump}]
        ]
        [Zero-state response $y_{zs}(t)$
            [determined by\\the input alone]
            [zero initial state]
        ]
        [Natural response
            [homogeneous-solution form]
            [coefficients fixed by initial state\\and input together]
        ]
        [Forced response
            [particular-solution form]
            [determined by the input]
        ]
    ]
\end{forest}}
\caption{Two decompositions of the complete response}
\end{figure}

\section{Impulse Response and Step Response}

\subsection{Impulse Response}

Define
\[
    h(t)=T\bigl[\{0\},\delta(t)\bigr].
\]

\textbf{How to find the impulse response}

Two cases arise for LTI systems: the right-hand side of the differential
equation contains only $f(t)$, or it contains $f(t)$ together with its
derivatives. Using the linearity and the differentiation property of the
zero-state response of an LTI system, the second case reduces to the first:
\[
    \sum_{j=0}^{n}a_j\,y^{(j)}(t)=\sum_{j=0}^{m}b_j\,f^{(j)}(t).
\]
Choose a new variable $y_1(t)$ whose differential equation has the same
left-hand side but only $f(t)$ on the right:
\[
    \sum_{j=0}^{n}a_j\,y_1^{(j)}(t)=f(t).
\]
Let its impulse response be $h_1(t)$. By linearity and the differentiation
property,
\[
    h(t)=b_m h_1^{(m)}(t)+b_{m-1}h_1^{(m-1)}(t)+\cdots+b_0 h_1(t).
\]

\subsection{Step Response}

Define
\[
    g(t)=T\bigl[\{0\},\varepsilon(t)\bigr].
\]
If the right-hand side of the equation contains only $f(t)$, then
\[
    \begin{cases}
        g^{(n)}(t)+a_{n-1}g^{(n-1)}(t)+\cdots+a_0g(t)=\varepsilon(t),\\[4pt]
        g^{(j)}(0_{-})=0,\quad j=0,1,2,\dots,n-1 .
    \end{cases}
\]
If the characteristic roots are all simple, then
\[
    g(t)=\Bigl(\sum_{j=1}^{n}C_j e^{\lambda_j t}+\frac{1}{a_0}\Bigr)\varepsilon(t).
\]
Since $\varepsilon(t)$ is the integral of $\delta(t)$, $g(t)$ is the integral
of $h(t)$: $g(t)=\displaystyle\int_{-\infty}^{t}h(\tau)\,\dd\tau$; finding one
of them gives the other.

\section{Zero-State Response and the Convolution Integral}

The zero-state response of an LTI system equals the convolution of the input
with the impulse response. Decompose the input into a superposition of
impulses,
\[
    f(t)=\int_{-\infty}^{\infty}f(\tau)\,\delta(t-\tau)\,\dd\tau ,
\]
By linearity and time invariance, each impulse $\delta(t-\tau)$ produces the
response $h(t-\tau)$, hence
\begin{equation}
    y_{zs}(t)=\int_{-\infty}^{\infty}f(\tau)\,h(t-\tau)\,\dd\tau
            =\int_{-\infty}^{\infty}h(\tau)\,f(t-\tau)\,\dd\tau
            \overset{\mathrm{def}}{=}f(t)*h(t).
\end{equation}
This is the \textbf{convolution integral}. It obeys the commutative,
distributive, and associative laws, together with convolution by an impulse,
$f(t)*\delta(t-t_0)=f(t-t_0)$; for a causal system $h(t)$ vanishes for $t<0$,
so the lower limit of the integral may be changed to $0$. With a causal input,
\[
    y_{zs}(t)=f(t)*h(t)=\int_{0}^{t}f(\tau)h(t-\tau)\,\dd\tau .
\]

\section{Zero-State Response and the Convolution Integral}
We defined a narrow pulse $p_n(t)$ of unit strength (i.e.\ unit area under the
pulse waveform).
\\
An arbitrary input $f(t)$ can be split into many narrow pulses of width
$\Delta \tau$.
\\
The $k$-th pulse occurs at time $t=k\Delta\tau$ with height $f(k\Delta\tau)$,
and can be written as $f(k\Delta \tau)\Delta \tau\, p_n(t-k\Delta \tau)$:
\begin{equation}
    f_{\Delta}(t)=\sum_{k=-\infty}^{\infty}f(k\Delta \tau)\Delta \tau p_n(t-k\Delta \tau)
\end{equation}
whose limit as $\Delta \tau \rightarrow 0$ is $f(t)$:
\begin{equation}
    f(t)=\lim_{\tau \rightarrow 0}f_{\Delta}(t)=\int_{-\infty}^{\infty}f(\tau)\delta(t-\tau)d\tau
\end{equation}
Through this process, $f(t)$ is microscopically split into a sum of
infinitely many impulse functions. By the linearity and time invariance of
LTI systems, the zero-state response to an arbitrary input is
\begin{equation}
    y_{zs}(t)=\int_{-\infty}^{\infty}f(\tau)h(t-\tau)d\tau
\end{equation}
This operation is defined as the convolution integral. The equation above
states that the zero-state response $y_{zs}(t)$ of an LTI system is the
convolution of $f(t)$ with the impulse response $h(t)$, written as
\begin{equation}
    y_{zs}(t)=f(t)* h(t)
\end{equation}
Microscopically, the input $f(t)$ is decomposed into a sum of infinitely many
impulses, each with strength equal to the value of $f(t)$ at its position; by
linearity, the response is the sum of the responses to all the impulses acting
separately.
By time invariance, each such response is time-shifted as well. By linearity,
the response of the system at $t=t_0$ equals the sum of the responses produced
by all input components acting separately.

\subsection{Properties of the Convolution Integral}
\begin{equation}
    f(t)=f_1(t)*f_2(t)=\int_{-\infty}^{\infty}f_1(\tau)f_2(t-\tau)d\tau
\end{equation}
\textbf{Property 1 (commutativity)}\\
\begin{equation}
    f_1(t)*f_2(t)=f_2(t)*f_1(t)
\end{equation}
\begin{note}
\textbf{Proof}
\begin{equation*}
    \int_{-\infty}^{\infty}f_1(\tau)f_2(t-\tau)d\tau = \int_{-\infty}^{\infty}f_2(\tau)f_1(t-\tau)d\tau
\end{equation*}
\end{note}
\textbf{Property 2 (distributivity)}\\
\begin{equation}
    f_1(t)*[f_2(t)+f_3(t)]=f_1(t)*f_2(t)+f_1(t)*f_3(t)
\end{equation}
\begin{note}
\textbf{Proof}
\begin{equation*}
    \int_{-\infty}^{\infty}f_1(\tau)[f_2(t-\tau)+f_3(t-\tau)]d\tau = \int_{-\infty}^{\infty}f_1(\tau)f_2(t-\tau)d\tau + \int_{-\infty}^{\infty}f_1(\tau)f_3(t-\tau)d\tau
\end{equation*}
\end{note}
\textbf{Property 3 (associativity)}\\
\begin{equation}
    f_1(t)*[f_2(t)*f_3(t)]=[f_1(t)*f_2(t)]*f_3(t)
\end{equation}
\begin{note}
\textbf{Proof}
\begin{equation*}
    \int_{-\infty}^{\infty}f_1(\tau)[\int_{-\infty}^{\infty}f_2(\sigma)f_3(t-\tau-\sigma)d\sigma]d\tau = \int_{-\infty}^{\infty}[\int_{-\infty}^{\infty}f_1(\tau)f_2(\sigma)d\tau]f_3(t-\sigma)d\sigma
\end{equation*}
\end{note}

% ==================================================================
\subsection{Convolution with the Impulse Function}
% ==================================================================
Using commutativity, the convolution of a function with the impulse function is
\[
f(t)*\delta(t)=\delta(t)*f(t)=\int_{-\infty}^{\infty}\delta(\tau)f(t-\tau)d\tau=f(t)
\]
that is,
\[
f(t)*\delta(t)=\delta(t)*f(t)=f(t)
\]
More generally,
\[
f(t)*\delta(t-t_0)=\delta(t-t_0)*f(t)=f(t-t_0)
\]
In addition,
\[
f(t-t_1)*\delta(t-t_2)=\delta(t-t_2)*f(t-t_1)=f(t-t_1-t_2)
\]

\subsection{Differentiation and Integration of Convolution}
\begin{definition*}
   \[ f^{(1)}(t)=\frac{df(t)}{dt}
   \]
\end{definition*}
\begin{definition*}
   \[ f^{(-1)}(t)=\int_{-\infty}^t f(\tau)d\tau
   \]
\end{definition*}
If
\[
f(t)=f_1(t)*f_2(t)=f_2(t)*f_1(t)
\]
then
\[
    f^{(1)}(t)=f_1^{(1)}(t)*f_2(t)=f_1(t)*f_2^{(1)}(t)
\]
\[
    f^{(-1)}(t)=f_1^{(-1)}(t)*f_2(t)=f_1(t)*f_2^{(-1)}(t)
\]
\textbf{Generalized result}
\[
    f^{(i)}(t)=f_1^{(j)}(t)*f_2^{(i-j)}(t)
\]
% ==================================================================
\chapter{Time-Domain Analysis of Discrete-Time Systems}
% ==================================================================

\section{Differences and Difference Equations}

Define the first-order forward difference as
\[
    \Delta f(k)=f(k+1)-f(k),
\]
and the first-order backward difference as
\[
    \nabla f(k)=f(k)-f(k-1).
\]
The two behave identically; we use the backward difference below. Verify
linearity:
\[
\begin{aligned}
\nabla\bigl[a_1f_1(k)+a_2f_2(k)\bigr]
&= \bigl[a_1f_1(k)+a_2f_2(k)\bigr]-\bigl[a_1f_1(k-1)+a_2f_2(k-1)\bigr]\\
&= a_1\bigl[f_1(k)-f_1(k-1)\bigr]+a_2\bigl[f_2(k)-f_2(k-1)\bigr]\\
&= a_1\nabla f_1(k)+a_2\nabla f_2(k).
\end{aligned}
\]
Hence the difference operator is linear.

Define the $n$-th order difference
\[
    \nabla^{n}f(k)=\nabla\bigl[\nabla^{n-1}f(k)\bigr]
    =\sum_{j=0}^{n}(-1)^{j}\dbinom{n}{j}f(k-j).
\]
A difference equation is an equation in the independent variable $k$ involving
an unknown sequence $y(k)$ and its differences, generally written
\[
    F\bigl(k,y(k),\nabla y(k),\dots,\nabla^{n}y(k)\bigr)=0 .
\]
Since an $n$-th order difference is a linear combination of $y(k)$ and its
shifted versions, it may equally be written
\[
    G\bigl(k,y(k),y(k-1),y(k-2),\dots,y(k-n)\bigr)=0 .
\]
Discrete LTI systems are described by linear difference equations with
constant coefficients, which can be solved numerically by iteration.

\section{The Classical Solution of Difference Equations}

For a single-input, single-output discrete LTI system with input $f(k)$ and
complete response $y(k)$, the model is an $n$-th order linear
constant-coefficient difference equation:
\[
    \sum_{j=0}^{n}a_{n-j}\,y(k-j)=\sum_{i=0}^{m}b_{m-i}\,f(k-i).
\]
Its solution again combines a homogeneous and a particular part:
\[
    y(k)=y_h(k)+y_p(k).
\]

\subsection{The Homogeneous Solution}

The solution of the homogeneous equation
\[
    y(k)+a_{n-1}y(k-1)+a_{n-2}y(k-2)+\cdots+a_0y(k-n)=0
\]
is called the homogeneous solution. Take a first-order difference equation as
an example:
\[
    y(k)+ay(k-1)=0
    \quad\Longrightarrow\quad
    \frac{y(k)}{y(k-1)}=-a,
\]
which is a geometric sequence, so $y(k)$ must have the form
\[
    y(k)=C_0(-a)^{k}.
\]
For an $n$-th order equation the homogeneous solution is a combination of
sequences $C\lambda^{k}$. Substituting into the homogeneous equation:
\[
    C\lambda^{k}+a_{n-1}C\lambda^{k-1}+a_{n-2}C\lambda^{k-2}+\cdots+a_0C\lambda^{k-n}=0 .
\]
Cancelling $C$ and $\lambda^{k-n}$:
\[
    \lambda^{n}+a_{n-1}\lambda^{n-1}+\cdots+a_1\lambda+a_0=0 .
\]
This is the \textbf{characteristic equation}; different roots give different
forms of the homogeneous solution:

\begin{table}[htbp]
\centering
\caption{Forms of the homogeneous solution $y_h(k)$}
\small
\begin{tabular}{@{}p{0.30\linewidth}p{0.62\linewidth}@{}}
\toprule
Characteristic root $\lambda$ & Homogeneous solution $y_h(k)$ \\
\midrule
simple real root & $C\lambda^{k}$ \\
real root of multiplicity $r$ & $\Bigl(\sum_{i=0}^{r-1}C_i k^{i}\Bigr)\lambda^{k}$ \\
a pair of complex conjugate roots $\lambda_{1,2}=a\pm bi=\rho e^{\pm j\theta}$ &
$\rho^{k}\bigl(C\cos\beta k+D\sin\beta k\bigr)$ or $A\rho^{k}\cos(\beta k+\theta)$ \\
complex conjugate roots of multiplicity $r$ &
$\rho^{k}\Bigl(\sum_{i=0}^{r-1}A_i k^{r-1}\cos(\beta k+\theta_i)\Bigr)$ \\
\bottomrule
\end{tabular}
\end{table}

\subsection{The Particular Solution}

The functional form of the particular solution depends on the form of the
input:

\begin{table}[htbp]
\centering
\caption{Forms of the particular solution $y_p(k)$}
\small
\begin{tabular}{@{}p{0.28\linewidth}p{0.42\linewidth}p{0.24\linewidth}@{}}
\toprule
Input $f(k)$ & Particular solution $y_p(k)$ & Condition \\
\midrule
$k^{m}$ & $P_mk^{m}+P_{m-1}k^{m-1}+\cdots+P_1k+P_0$ & no characteristic root equals $1$ \\
$k^{m}$ & $k^{r}\bigl(P_mk^{m}+P_{m-1}k^{m-1}+\cdots+P_1k+P_0\bigr)$ & characteristic root $1$ of multiplicity $r$ \\
$a^{k}$ & $Pa^{k}$ & $a$ is not a characteristic root \\
$a^{k}$ & $\bigl(P_rk^{r}+P_{r-1}k^{r-1}+\cdots+P_1k+P_0\bigr)a^{k}$ & $a$ is a characteristic root of multiplicity $r$ \\
$\cos\beta k$ or $\sin\beta k$ & $P\cos\beta k+Q\sin\beta k$ & no characteristic root equals $e^{j\beta}$ \\
\bottomrule
\end{tabular}
\end{table}

\subsection{Zero-Input and Zero-State Responses}

The response caused by the initial state alone (zero input) is the zero-input
response:
\[
    \sum_{j=0}^{n}a_{n-j}\,y_{zi}(k-j)=0 .
\]
The response of a zero-initial-state system to the input alone is the
zero-state response:
\[
    \begin{cases}
        \displaystyle\sum_{j=0}^{n}a_{n-j}\,y_{zs}(k-j)=\sum_{i=0}^{m}b_{m-i}\,f(k-i),\\[8pt]
        y_{zs}(-1)=y_{zs}(-2)=\cdots=y_{zs}(-n)=0 .
    \end{cases}
\]

\begin{note}
The complete response can be decomposed into the zero-input and zero-state
responses, or into the natural and forced responses; the two decompositions
differ slightly. Both the zero-input response and the natural response are
homogeneous solutions, but with different coefficients: the coefficients of
the natural response are fixed by the initial state \emph{and} the input
together, whereas those of the zero-input response are fixed by the initial
state alone.
\end{note}

\section{The Unit Sequence and the Unit Sequence Response}

The unit sequence is defined as
\[
    \delta(k)=
    \begin{cases}
        1, & k=0,\\
        0, & k\neq 0 .
    \end{cases}
\]
The unit step sequence is defined as
\[
    \varepsilon(k)=
    \begin{cases}
        0, & k<0,\\
        1, & k\ge 0 .
    \end{cases}
\]
They are related by $\delta(k)=\varepsilon(k)-\varepsilon(k-1)$ and
$\varepsilon(k)=\sum_{i=0}^{\infty}\delta(k-i)$.

\subsection{The Unit Sequence Response}

The unit sequence response is denoted $h(k)$. It can be found by solving the
difference equation, or by the $z$-transform. When the right-hand side of the
equation contains shifted copies of $f(k)$, the linearity and shift invariance
of the zero-state response reduce the problem to the case where the
right-hand side contains only $f(k)$.

\section{Zero-State Response and the Convolution Sum}

\subsection{Zero-State Response to an Arbitrary Input}

For an LTI discrete system, decompose the input into unit sequences: if the
unit sequence response is known, the response to each unit sequence follows
immediately, and their sum is the zero-state response. This summation is
called the \textbf{convolution sum}.
\[
    f(k)=\sum_{i=-\infty}^{\infty}f(i)\,\delta(k-i),
\]
hence
\[
    y_{zs}(k)=\sum_{i=-\infty}^{\infty}f(i)\,h(k-i)=f(k)*h(k).
\]

\subsection{The Convolution Sum}

\begin{equation}
    y(k)=f_1(k)*f_2(k)=\sum_{i=-\infty}^{\infty}f_1(i)\,f_2(k-i).
\end{equation}
For causal sequences (zero for $k<0$) the lower limit becomes $0$ and the
upper limit $k$:
\[
    y(k)=\sum_{i=0}^{k}f_1(i)\,f_2(k-i).
\]

\subsection{Properties of the Convolution Sum}

\begin{itemize}
    \item \textbf{Commutativity}: $f_1(k)*f_2(k)=f_2(k)*f_1(k)$;
    \item \textbf{Distributivity}: $f_1(k)*\bigl[f_2(k)+f_3(k)\bigr]=f_1(k)*f_2(k)+f_1(k)*f_3(k)$;
    \item \textbf{Associativity}: $\bigl[f_1(k)*f_2(k)\bigr]*f_3(k)=f_1(k)*\bigl[f_2(k)*f_3(k)\bigr]$;
    \item \textbf{Time shifting}: if $f_1(k)*f_2(k)=f(k)$, then
          $f_1(k-k_1)*f_2(k-k_2)=f(k-k_1-k_2)$;
    \item \textbf{Convolution with the unit sequence}: $f(k)*\delta(k)=f(k)$,
          $f(k)*\delta(k-k_0)=f(k-k_0)$;
    \item \textbf{Accumulation}: $\displaystyle\sum_{i=-\infty}^{k}f(i)=f(k)*\varepsilon(k)$.
\end{itemize}

% ==================================================================
\chapter{The Fourier Transform and Frequency-Domain Analysis}
% ==================================================================

\section{Notational Conventions}

\begin{table}[htbp]
\centering
\caption{Notation used in this chapter}
\begin{tabular}{@{}ll@{}}
\toprule
Symbol & Meaning \\
\midrule
$T$ & period of the signal \\
$\Omega=\dfrac{2\pi}{T}$ & fundamental angular frequency (periodic signals) \\
$n$ & harmonic order, $n=0,1,2,\dots$ \\
$\omega$ & continuous angular frequency (aperiodic signals) \\
$\tau$ & rectangular pulse width; also a dummy integration variable \\
\bottomrule
\end{tabular}
\end{table}

\section{Orthogonal Functions and Orthogonal Decomposition of Signals}

\subsection{Sets of Orthogonal Functions}

\begin{definition*}
If
\[
    \int_{t_1}^{t_2}\varphi_1(t)\,\varphi_2(t)\,\dd t=0,
\]
then $\varphi_1(t)$ and $\varphi_2(t)$ are said to be \textbf{orthogonal} on
$[t_1,t_2]$. $n$ functions that are pairwise orthogonal form an
\textbf{orthogonal set} $\{\varphi_1,\varphi_2,\dots,\varphi_n\}$.
\end{definition*}

\textbf{Completeness}\quad If there is no function
$\phi(t)\notin\{\varphi_i\}$ such that for every $i$
\[
    \int_{t_1}^{t_2}\phi(t)\,\varphi_i(t)\,\dd t=0,
\]
then the set is a \textbf{complete orthogonal set}.

\textbf{Complex function sets}\quad The orthogonality condition becomes an
inner product with conjugation:
\[
    \int_{t_1}^{t_2}\varphi_i(t)\,\varphi_j^{*}(t)\,\dd t=
    \begin{cases}
        0, & i\neq j,\\[2pt]
        k_i\neq 0, & i=j .
    \end{cases}
\]

\begin{note}
\textbf{Why this analogy is legitimate}

Anything closed under linear operations can be regarded as a ``vector''. A
continuous function is an infinite-dimensional vector: its value at each point
of the independent variable is one coordinate,
\[
    f(x)=\lim_{\Delta x\to 0}\bigl(f(a),\,f(a+\Delta x),\,f(a+2\Delta x),\,\dots\bigr).
\]
Multiplying two functions pointwise on $[a,b]$ and summing (integrating) gives
their inner product; a vanishing inner product means ``perpendicular'' --- this
is exactly the geometric meaning of orthogonality.
\end{note}

\subsection{Decomposing a Signal into Orthogonal Functions}

Approximate with finitely many terms:
\[
    f(t)\approx\sum_{j=1}^{n}c_j\varphi_j(t).
\]
Measure the quality of the approximation by the \textbf{mean-square error}
(called the root-mean-square error in some texts):
\[
    \overline{\epsilon^{2}}\ \overset{\mathrm{def}}{=}\
    \frac{1}{t_2-t_1}\int_{t_1}^{t_2}
    \Bigl[f(t)-\sum_{j=1}^{n}c_j\varphi_j(t)\Bigr]^{2}\dd t .
\]
The error is squared so that positive and negative deviations cannot cancel
each other.

\textbf{Optimal coefficients}\quad Set
$\partial\overline{\epsilon^{2}}/\partial c_j=0$ and use orthogonality,
$\int\varphi_i\varphi_j=0\ (i\neq j)$:
\[
    \frac{\partial}{\partial c_j}\int_{t_1}^{t_2}
    \Bigl[f(t)-\sum_{i}c_i\varphi_i(t)\Bigr]^{2}\dd t
    =\int_{t_1}^{t_2}\bigl[-2f(t)\varphi_j(t)+2c_j\varphi_j^{2}(t)\bigr]\dd t=0,
\]
hence
\[
    \boxed{\;
    c_j=\frac{\displaystyle\int_{t_1}^{t_2}f(t)\varphi_j(t)\,\dd t}
             {\displaystyle\int_{t_1}^{t_2}\varphi_j^{2}(t)\,\dd t}
       =\frac{1}{k_j}\int_{t_1}^{t_2}f(t)\varphi_j(t)\,\dd t\;}
\]
where $k_j=\displaystyle\int_{t_1}^{t_2}\varphi_j^{2}(t)\,\dd t$.

\textbf{Minimum mean-square error}\quad Substituting back,
\[
    \overline{\epsilon^{2}}=\frac{1}{t_2-t_1}
    \left[\int_{t_1}^{t_2}f^{2}(t)\,\dd t-\sum_{j=1}^{n}c_j^{2}k_j\right].
\]
Since $\overline{\epsilon^{2}}\ge 0$, more terms mean smaller error; as
$n\to\infty$ with a complete set, $\overline{\epsilon^{2}}\to 0$, and then
\[
    \int_{t_1}^{t_2}f^{2}(t)\,\dd t=\sum_{j=1}^{\infty}c_j^{2}k_j,
    \qquad
    f(t)=\sum_{j=1}^{\infty}c_j\varphi_j(t).
\]
This energy equality is \textbf{Parseval's equation}.

\section{Fourier Series}

The Fourier series is Parseval's equation specialized to the trigonometric
orthogonal basis.

\subsection{Trigonometric Form}

\begin{equation}
    f(t)=\frac{a_0}{2}+\sum_{n=1}^{\infty}
    \bigl[a_n\cos(n\Omega t)+b_n\sin(n\Omega t)\bigr].
\end{equation}
The coefficients (from the general formula of the previous section with
$\varphi_j$ taken as $\cos n\Omega t$, $\sin n\Omega t$ and $k=\frac{T}{2}$):
\[
    a_n=\frac{2}{T}\int_{-\frac{T}{2}}^{\frac{T}{2}}f(t)\cos(n\Omega t)\,\dd t
    \quad(n\ge 0),
\]
\[
    b_n=\frac{2}{T}\int_{-\frac{T}{2}}^{\frac{T}{2}}f(t)\sin(n\Omega t)\,\dd t
    \quad(n\ge 1).
\]
Here $\dfrac{a_0}{2}=\dfrac{1}{T}\displaystyle\int_{-\frac{T}{2}}^{\frac{T}{2}}f(t)\,\dd t$
is the DC component.

\subsection{Amplitude--Phase Form}

Using the auxiliary-angle identity
$a_n\cos\theta+b_n\sin\theta=\sqrt{a_n^{2}+b_n^{2}}\,\cos\bigl(\theta-\arctan\frac{b_n}{a_n}\bigr)$:
\[
    f(t)=\frac{A_0}{2}+\sum_{n=1}^{\infty}A_n\cos(n\Omega t+\varphi_n),
\]
\[
    A_0=a_0,\qquad A_n=\sqrt{a_n^{2}+b_n^{2}},\qquad
    \varphi_n=-\arctan\frac{b_n}{a_n}.
\]
The term with $\omega=\Omega$ is the \textbf{fundamental}; terms at $n\Omega$
are the \textbf{$n$-th harmonics}.

\subsection{Convergence and the Gibbs Phenomenon}

More harmonics give a waveform closer to the original signal (except at
discontinuities). Near a discontinuity the spike shrinks toward the jump as
$n$ grows, but the \textbf{overshoot does not decrease with $n$}; its limit is
about $9\%$ --- the \textbf{Gibbs phenomenon}.

\begin{note}
\textbf{Derivation of the $9\%$ Gibbs overshoot}

Take the square wave
\[
    f(x)=
    \begin{cases}
        \ \ 1, & 0<x<\pi,\\
        -1, & -\pi<x<0,
    \end{cases}
    \qquad
    f(x)=\frac{4}{\pi}\sum_{k=0}^{\infty}\frac{\sin(2k+1)x}{2k+1}.
\]
The partial sum is
\[
    S_N(x)=\frac{4}{\pi}\sum_{k=0}^{N}\frac{\sin(2k+1)x}{2k+1}
    =\frac{4}{\pi}\int_0^{x}\sum_{k=0}^{N}\cos(2k+1)t\,\dd t .
\]
With the kernel identity
$\displaystyle\sum_{k=0}^{N}\cos(2k+1)t=\frac{\sin(2N+2)t}{2\sin t}$,
\[
    S_N(x)=\frac{4}{\pi}\int_0^{x}\frac{\sin(2N+2)t}{2\sin t}\,\dd t
    =\frac{2}{\pi}\int_0^{(2N+2)x}
    \frac{\sin u}{(2N+2)\sin\frac{u}{2N+2}}\,\dd u,
\]
where we substituted $u=(2N+2)t$, $\dd t=\dfrac{\dd u}{2N+2}$.

Let $N\to\infty$; then $(2N+2)\sin\frac{u}{2N+2}\to u$, and
\[
    \lim_{N\to\infty}S_N(x)=\frac{2}{\pi}\int_0^{\infty}\frac{\sin u}{u}\,\dd u
    \ \xrightarrow{\ X\to\infty\ }\
    F(X)=\frac{2}{\pi}\int_0^{X}\frac{\sin u}{u}\,\dd u .
\]
$F'(X)=0$ at $X=n\pi$, so the first extremum is at $X=\pi$:
\[
    F(\pi)=\frac{2}{\pi}\Si(\pi)=1.17898 .
\]
Relative to the jump of size $2$, the overshoot ratio is
\[
    \delta=\frac{1.17898-1}{2}\approx 0.0895\approx 9\% .
\]
\end{note}

\subsection{Even/Odd Symmetry and Half-Wave Symmetry}

\textbf{(1) $f(t)$ even} $\Rightarrow b_n=0$,
\[
    a_n=\frac{4}{T}\int_0^{\frac{T}{2}}f(t)\cos(n\Omega t)\,\dd t,\qquad
    A_n=|a_n|,\quad \varphi_n=m\pi .
\]

\textbf{(2) $f(t)$ odd} $\Rightarrow a_n=0$,
\[
    b_n=\frac{4}{T}\int_0^{\frac{T}{2}}f(t)\sin(n\Omega t)\,\dd t .
\]

\textbf{(3) Odd-harmonic functions (half-wave symmetry)}\quad If
\[
    f(t)=-f\!\left(t\pm\frac{T}{2}\right),
\]
then only \textbf{odd} harmonics appear; even harmonics are absent.

\subsection{Exponential Form}

\begin{equation}
    f(t)=\sum_{n=-\infty}^{\infty}F_n e^{jn\Omega t},
    \qquad
    F_n=\frac{1}{T}\int_{-\frac{T}{2}}^{\frac{T}{2}}f(t)e^{-jn\Omega t}\,\dd t .
\end{equation}
The complex coefficients can be written $F_n=|F_n|e^{j\varphi_n}$.

\begin{note}
\textbf{Key point of the derivation}\quad The self-product integral of the
basis $e^{jn\Omega t}$ is
$\int_{-T/2}^{T/2}e^{jn\Omega t}\bigl(e^{jn\Omega t}\bigr)^{*}\dd t=T$,
while distinct $n$ give $0$; thus $k=T$, and the general formula of the
previous section applies directly.
\end{note}

\subsection{Relation Between the Two Forms}

From Euler's formula
\[
    \cos x=\frac{e^{jx}+e^{-jx}}{2},
\]
\[
\begin{aligned}
    f(t)&=\frac{A_0}{2}+\sum_{n=1}^{\infty}\frac{A_n}{2}
          \Bigl[e^{j(n\Omega t+\varphi_n)}+e^{-j(n\Omega t+\varphi_n)}\Bigr]\\
        &=\frac{1}{2}\sum_{n=-\infty}^{\infty}A_n e^{j\varphi_n}e^{jn\Omega t}.
\end{aligned}
\]
Hence
\[
    F_n=\frac{1}{2}A_n e^{j\varphi_n},\qquad |F_n|=\frac{1}{2}A_n .
\]
That is, \textbf{the double-sided spectrum has half the amplitude of the
single-sided one}; the negative- and positive-frequency terms of the
double-sided spectrum together synthesize each physical component of the
single-sided spectrum.

\section{Spectra of Periodic Signals}

From $f(t)=\dfrac{A_0}{2}+\sum A_n\cos(n\Omega t+\varphi_n)$ or
$f(t)=\sum F_ne^{jn\Omega t}$, a periodic signal is completely determined by
$\Omega$, the amplitudes, and the initial phases of its harmonics. Plotting
amplitude or phase against $n\Omega$ (or $n$) gives the \textbf{spectrum}.

\begin{table}[htbp]
\centering
\caption{The two spectrum representations}
\begin{tabular}{@{}lll@{}}
\toprule
Form & Range of $n$ & Name \\
\midrule
trigonometric, $A_n,\varphi_n$ & $n\in[1,+\infty)$ & single-sided spectrum \\
exponential, $F_n$ & $n\in(-\infty,+\infty)$ & double-sided spectrum \\
\bottomrule
\end{tabular}
\end{table}

For \textbf{real signals}: the double-sided amplitude spectrum is even and the
double-sided phase spectrum is odd (this fails for complex signals).

\subsection{Spectrum of a Periodic Rectangular Pulse}

\[
\begin{aligned}
    F_n&=\frac{1}{T}\int_{-\frac{T}{2}}^{\frac{T}{2}}f(t)e^{-jn\Omega t}\,\dd t
        =\frac{1}{T}\int_{-\frac{\tau}{2}}^{\frac{\tau}{2}}e^{-jn\Omega t}\,\dd t\\
       &=\frac{2}{T}\cdot\frac{\sin\left(\dfrac{n\Omega\tau}{2}\right)}{n\Omega}
        =\frac{\tau}{T}\cdot
         \frac{\sin\left(\dfrac{n\Omega\tau}{2}\right)}{\dfrac{n\Omega\tau}{2}}
        =\frac{\tau}{T}\,\Sa\!\left(\frac{n\Omega\tau}{2}\right).
\end{aligned}
\]
where $\Sa(x)=\dfrac{\sin x}{x}$ is the sampling function.

\begin{itemize}
    \item The spectrum is \textbf{discrete}, with line spacing
          $\Omega=\dfrac{2\pi}{T}$: the larger $T$, the denser the lines;
    \item the envelope is governed by $\Sa$, with zeros at
          $\dfrac{n\Omega\tau}{2}=k\pi$, i.e.\ $f=\dfrac{k}{\tau}$;
    \item the energy is concentrated inside the first zero; define the
          \textbf{bandwidth}
    \[
        0\le f\le\frac{1}{\tau},\qquad \Delta F=\frac{1}{\tau}.
    \]
\end{itemize}

As $T\to\infty$ the line spacing tends to zero and the discrete spectrum
becomes continuous --- this is precisely how the Fourier transform arises.

\begin{figure}[htbp]
\centering
\caption{Spectrum of a periodic rectangular pulse (discrete lines with a
$\Sa$-shaped envelope)}
\end{figure}

\section{Power of Periodic Signals and Parseval's Identity}

The average power is
\[
    P=\frac{1}{T}\int_{-\frac{T}{2}}^{\frac{T}{2}}f^{2}(t)\,\dd t .
\]
Substituting the trigonometric series and using orthogonality of the
trigonometric system:
\[
    P=\frac{A_0^{2}}{4}+\sum_{n=1}^{\infty}\frac{A_n^{2}}{2}
     =P_{\text{DC}}+P_{\text{AC}} .
\]
Since $|F_n|=\frac{1}{2}A_n$,
\[
    \boxed{\;P=\sum_{n=-\infty}^{\infty}|F_n|^{2}\;}
\]
--- \textbf{Parseval's identity}: the time-domain average power equals the sum
of the powers of all spectral lines.

\section{The Fourier Transform}

\subsection{From Series to Transform}

A periodic signal $f_T(t)=\sum F_ne^{jn\Omega t}$ with
$F_n=\frac{1}{T}\int_{-T/2}^{T/2}f_T(t)e^{-jn\Omega t}\dd t$. Let $T\to\infty$:

\begin{itemize}
    \item the line spacing $\Omega\to\dd\omega$, and the discrete sum becomes
          an integral;
    \item the amplitudes tend to zero, so the \textbf{spectral density}
          $F(j\omega)$ replaces $F_n$.
\end{itemize}

\[
    f(t)=\frac{1}{2\pi}\int_{-\infty}^{\infty}
    \left[\int_{-\infty}^{\infty}f(\tau)e^{-j\omega\tau}\,\dd\tau\right]
    e^{j\omega t}\,\dd\omega .
\]
Hence the definitions
\[
    \boxed{\;F(j\omega)=\int_{-\infty}^{\infty}f(t)e^{-j\omega t}\,\dd t\;}
    \qquad\text{(Fourier transform)}
\]
\[
    \boxed{\;f(t)=\frac{1}{2\pi}\int_{-\infty}^{\infty}F(j\omega)e^{j\omega t}\,\dd\omega\;}
    \qquad\text{(inverse Fourier transform)}
\]
written $\mathcal{F}[f(t)]=F(j\omega)$,
$\mathcal{F}^{-1}[F(j\omega)]=f(t)$. $F(j\omega)$ is called the
\textbf{spectral density function}.

\subsection{Two Ways of Writing the Complex Spectrum}

\[
    F(j\omega)=|F(j\omega)|e^{j\varphi(\omega)}=R(\omega)+jX(\omega).
\]
From the inverse transform (using that $|F|$ is even and $\varphi$ is odd),
the real and imaginary parts correspond to cosine and sine components:
\[
    f(t)=\frac{1}{\pi}\int_0^{\infty}|F(j\omega)|\cos\bigl(\omega t+\varphi(\omega)\bigr)\,\dd\omega .
\]
Thus a real aperiodic signal is a continuous superposition of cosines of
amplitude $\dfrac{|F(j\omega)|\dd\omega}{\pi}$.

\subsection{Even/Odd Symmetry of Signals and Symmetry of the Spectrum}

Substituting $e^{-j\omega t}=\cos\omega t-j\sin\omega t$ into the definition:
\[
    F(j\omega)=\underbrace{\int_{-\infty}^{\infty}f(t)\cos\omega t\,\dd t}_{R(\omega)}
    -j\underbrace{\int_{-\infty}^{\infty}f(t)\sin\omega t\,\dd t}_{-X(\omega)},
\]
\[
    R(\omega)=\int_{-\infty}^{\infty}f(t)\cos\omega t\,\dd t
    \quad(\text{even in }\omega),
\]
\[
    X(\omega)=-\int_{-\infty}^{\infty}f(t)\sin\omega t\,\dd t
    \quad(\text{odd in }\omega).
\]
Therefore
\[
    |F(-j\omega)|=|F(j\omega)|\ \ (\text{even}),\qquad
    \varphi(-\omega)=-\varphi(\omega)\ \ (\text{odd}),
\]
\[
    \boxed{\;F(-j\omega)=F^{*}(j\omega)\;}\qquad(\text{real signal, conjugate symmetry}).
\]
The phase angle is
\[
    \varphi(\omega)=\arctan\frac{X(\omega)}{R(\omega)} .
\]

\begin{table}[htbp]
\centering
\caption{Special cases for real signals}
\small
\begin{tabular}{@{}p{0.20\linewidth}p{0.74\linewidth}@{}}
\toprule
$f(t)$ & Conclusion \\
\midrule
real \textbf{even} function & $X(\omega)=0$,
$F(j\omega)=R(\omega)=2\displaystyle\int_0^{\infty}f(t)\cos\omega t\,\dd t$,
real and even \\
real \textbf{odd} function & $R(\omega)=0$,
$F(j\omega)=jX(\omega)=-2j\displaystyle\int_0^{\infty}f(t)\sin\omega t\,\dd t$,
purely imaginary and odd \\
\bottomrule
\end{tabular}
\end{table}

\textbf{Time reversal}\quad Let $\tau=-t$; then
\[
    \mathcal{F}[f(-t)]=\int_{-\infty}^{\infty}f(-t)e^{-j\omega t}\,\dd t
    =\int_{-\infty}^{\infty}f(\tau)e^{j\omega\tau}\,\dd\tau
    =F(-j\omega)=F^{*}(j\omega).
\]

\section{Fourier Transforms of Singularity Functions}

\begin{note}
\textbf{Remark}\quad The following signals are not absolutely integrable (one
of the Dirichlet conditions); strictly speaking the results are limits, or
transforms in the sense of generalized functions.
\end{note}

\textbf{(1) Unit impulse} $\delta(t)$
\[
    \mathcal{F}[\delta(t)]=\int_{-\infty}^{\infty}\delta(t)e^{-j\omega t}\,\dd t
    =e^{0}=1 .
\]
Most concentrated in time $\Rightarrow$ perfectly uniform in frequency (a
constant spectrum).

\textbf{(2) Derivatives of the impulse}\quad Using
$\displaystyle\int\delta^{(n)}(t)\varphi(t)\dd t=(-1)^{n}\varphi^{(n)}(0)$:
\[
    \mathcal{F}[\delta^{(n)}(t)]
    =\int_{-\infty}^{\infty}\delta^{(n)}(t)e^{-j\omega t}\,\dd t
    =(-1)^{n}\left.\frac{\dd^{n}}{\dd t^{n}}e^{-j\omega t}\right|_{t=0}
    =(j\omega)^{n}.
\]

\textbf{(3) Unit DC} $f(t)=1$

Take $f_1(t)=e^{-\alpha|t|}\xrightarrow{\alpha\to0}1$, while
\[
    F_1(j\omega)=\frac{2\alpha}{\alpha^{2}+\omega^{2}}
    \ \xrightarrow{\alpha\to0}\
    \begin{cases}
        \infty, & \omega=0,\\
        0, & \omega\neq 0,
    \end{cases}
\]
with total strength
\[
    \lim_{\alpha\to0}\int_{-\infty}^{\infty}\frac{2\alpha}{\alpha^{2}+\omega^{2}}\,\dd\omega
    =2\arctan\frac{\omega}{\alpha}\Big|_{-\infty}^{+\infty}=2\pi .
\]
Therefore
\[
    \mathcal{F}[1]=2\pi\delta(\omega).
\]

\textbf{(4) Signum function} $\sgn(t)$

Take $f_2(t)=\sgn(t)e^{-\alpha|t|}$, so
$F_2(j\omega)=\dfrac{-2j\omega}{\alpha^{2}+\omega^{2}}$; letting $\alpha\to0$:
\[
    \mathcal{F}[\sgn(t)]=\frac{2}{j\omega}\qquad(\omega\neq 0).
\]

\textbf{(5) Unit step}\quad
$\varepsilon(t)=\dfrac{1}{2}+\dfrac{1}{2}\sgn(t)$
\[
    \mathcal{F}[\varepsilon(t)]=\pi\delta(\omega)+\frac{1}{j\omega}
    =\pi\delta(\omega)-j\frac{1}{\omega}.
\]

\section{Properties of the Fourier Transform}

\begin{table}[htbp]
\centering
\caption{Main properties of the Fourier transform
($f(t)\lrarr F(j\omega)$)}
\small
\begin{tabular}{@{}p{0.26\linewidth}p{0.34\linewidth}p{0.34\linewidth}@{}}
\toprule
Property & Time domain & Frequency domain \\
\midrule
linearity & $af_1(t)+bf_2(t)$ & $aF_1(j\omega)+bF_2(j\omega)$ \\
duality & $F(jt)$ & $2\pi f(-\omega)$ \\
scaling & $f(at)$ & $\dfrac{1}{|a|}F\!\left(j\dfrac{\omega}{a}\right)$ \\
time shifting & $f(t\pm t_0)$ & $e^{\pm j\omega t_0}F(j\omega)$ \\
frequency shifting & $f(t)e^{\pm j\omega_0 t}$ & $F\bigl(j(\omega\mp\omega_0)\bigr)$ \\
time-domain convolution & $f_1(t)*f_2(t)$ & $F_1(j\omega)F_2(j\omega)$ \\
frequency-domain convolution & $f_1(t)f_2(t)$ & $\dfrac{1}{2\pi}F_1(j\omega)*F_2(j\omega)$ \\
time-domain differentiation & $f^{(n)}(t)$ & $(j\omega)^{n}F(j\omega)$ \\
time-domain integration & $f^{(-1)}(t)=\displaystyle\int_{-\infty}^{t}f(x)\dd x$
         & $\pi F(0)\delta(\omega)+\dfrac{F(j\omega)}{j\omega}$ \\
frequency-domain differentiation & $(-jt)^{n}f(t)$ & $F^{(n)}(j\omega)$ \\
frequency-domain integration & $\pi f(0)\delta(t)+\dfrac{f(t)}{-jt}$ & $F^{(-1)}(j\omega)$ \\
\bottomrule
\end{tabular}
\end{table}

\subsection{Linearity}

If
\[
    f_1(t)\lrarr F_1(j\omega),\qquad f_2(t)\lrarr F_2(j\omega),
\]
then
\[
    af_1(t)+bf_2(t)\lrarr aF_1(j\omega)+bF_2(j\omega).
\]

\subsection{Duality}

If $f(t)\lrarr F(j\omega)$, then
\[
    F(jt)\lrarr 2\pi f(-\omega).
\]

\begin{note}
\textbf{Proof}\quad From the inverse transform
\[
    f(t)=\frac{1}{2\pi}\int_{-\infty}^{\infty}F(j\omega)e^{j\omega t}\,\dd\omega,
\]
\[
    f(-t)=\frac{1}{2\pi}\int_{-\infty}^{\infty}F(j\omega)e^{-j\omega t}\,\dd\omega .
\]
Swapping the names of the variables $t\leftrightarrow\omega$,
\[
    f(-\omega)=\frac{1}{2\pi}\int_{-\infty}^{\infty}F(jt)e^{-j\omega t}\,\dd t
    =\frac{1}{2\pi}\mathcal{F}[F(jt)],
\]
so $F(jt)\lrarr 2\pi f(-\omega)$.
\end{note}

\subsection{Scaling}

\[
    f(at)\lrarr \frac{1}{|a|}F\!\left(j\frac{\omega}{a}\right).
\]

\begin{note}
\textbf{Proof}\quad
$\mathcal{F}[f(at)]=\displaystyle\int_{-\infty}^{\infty}f(at)e^{-j\omega t}\,\dd t$.
For $a>0$, set $at=x$; then
\[
    \mathcal{F}[f(at)]=\frac{1}{a}\int_{-\infty}^{\infty}f(x)e^{-j\omega\frac{x}{a}}\,\dd x
    =\frac{1}{a}F\!\left(j\frac{\omega}{a}\right).
\]
For $a<0$ the limits of integration reverse and an extra minus sign appears;
combining it with $|a|$ gives the general formula.
\end{note}

The duration of a signal and its occupied bandwidth are inversely related.

\subsection{Time Shifting}

If $f(t)\lrarr F(j\omega)$ and $t_0$ is a constant, then
\[
    f(t\pm t_0)\lrarr e^{\pm j\omega t_0}F(j\omega).
\]

\begin{note}
\textbf{Proof}\quad Let $x=t\pm t_0$:
\[
    \int_{-\infty}^{\infty}f(t\pm t_0)e^{-j\omega t}\,\dd t
    =\int_{-\infty}^{\infty}f(x)e^{-j\omega(x\mp t_0)}\,\dd x
    =e^{\pm j\omega t_0}F(j\omega).
\]
\end{note}

\subsection{Frequency Shifting}

If $f(t)\lrarr F(j\omega)$, then
\[
    f(t)e^{\pm j\omega_0 t}\lrarr F\bigl(j(\omega\mp\omega_0)\bigr).
\]
This is the theoretical basis of modulation (spectrum shifting).

\subsection{The Convolution Theorem}

\textbf{Time-domain convolution theorem}\quad If
$f_1(t)\lrarr F_1(j\omega)$ and $f_2(t)\lrarr F_2(j\omega)$, then
\[
    \mathcal{F}[f_1(t)*f_2(t)]=F_1(j\omega)F_2(j\omega).
\]

\begin{note}
\textbf{Proof}\quad
\[
\begin{aligned}
    \mathcal{F}[f_1(t)*f_2(t)]
    &=\int_{-\infty}^{\infty}\left[\int_{-\infty}^{\infty}f_1(\tau)f_2(t-\tau)\,\dd\tau\right]
      e^{-j\omega t}\,\dd t\\
    &=\int_{-\infty}^{\infty}f_1(\tau)
      \left[\int_{-\infty}^{\infty}f_2(t-\tau)e^{-j\omega t}\,\dd t\right]\dd\tau\\
    &=\int_{-\infty}^{\infty}f_1(\tau)e^{-j\omega\tau}F_2(j\omega)\,\dd\tau
      =F_1(j\omega)F_2(j\omega).
\end{aligned}
\]
\end{note}

\textbf{Frequency-domain convolution theorem}\quad If
$f_1(t)\lrarr F_1(j\omega)$ and $f_2(t)\lrarr F_2(j\omega)$, then
\[
    f_1(t)f_2(t)\lrarr\frac{1}{2\pi}F_1(j\omega)*F_2(j\omega).
\]
The proof is analogous: use the definitions, interchange the order of
integration, and apply the frequency-shifting property.

\subsection{Time-Domain Differentiation and Integration}

Define
\[
    f^{(n)}(t)=\frac{\dd^{n}f(t)}{\dd t^{n}},\qquad
    f^{(-1)}(t)=\int_{-\infty}^{t}f(x)\,\dd x .
\]
If $f(t)\lrarr F(j\omega)$, then
\[
    f^{(n)}(t)\lrarr(j\omega)^{n}F(j\omega),
\]
\[
    f^{(-1)}(t)\lrarr\pi F(0)\delta(\omega)+\frac{F(j\omega)}{j\omega},
\]
where
\[
    F(0)=\int_{-\infty}^{\infty}f(t)\,\dd t .
\]

\begin{note}
Sometimes $f(t)=g'(t)$ but
$g(t)\neq f^{(-1)}(t)=\displaystyle\int_{-\infty}^{t}f(x)\dd x$. Whether the
two coincide depends on whether $g(-\infty)$ is zero: only when
$g(-\infty)=0$ does $g(t)$ equal the running integral of $f(t)$; otherwise
they differ by a constant (which appears in the frequency domain as the term
$\pi F(0)\delta(\omega)$).
\end{note}

\subsection{Frequency-Domain Differentiation and Integration}

If $f(t)\lrarr F(j\omega)$, then
\[
    (-jt)^{n}f(t)\lrarr F^{(n)}(j\omega),
\]
\[
    \pi f(0)\delta(t)+\frac{f(t)}{-jt}\lrarr F^{(-1)}(j\omega),
\]
where
\[
    f(0)=\frac{1}{2\pi}\int_{-\infty}^{\infty}F(j\omega)\,\dd\omega .
\]

\section{Fourier Transform of Periodic Signals}

For the periodic signal
\[
    f_T(t)=\sum_{n=-\infty}^{\infty}F_ne^{jn\Omega t},
    \qquad
    F_n=\frac{1}{T}\int_{-\frac{T}{2}}^{\frac{T}{2}}f(t)e^{-jn\Omega t}\,\dd t,
\]
we get
\[
    \mathcal{F}[f_T(t)]=2\pi\sum_{n=-\infty}^{\infty}F_n\,\delta(\omega-n\Omega),
\]
i.e.\ the spectrum of a periodic signal is a train of equally spaced impulses
of strength $2\pi F_n$.

\subsection{Relation Between Fourier Coefficients and the Transform}

Extending a single pulse $f_0(t)$ periodically with period $T$ can be written
as a convolution with a periodic impulse train:
\[
    f_T(t)=f_0(t)*\delta_T(t),
    \qquad
    \delta_T(t)=\sum_{n=-\infty}^{\infty}\delta(t-nT).
\]
By the time-domain convolution theorem and
$\mathcal{F}[\delta_T(t)]=\Omega\displaystyle\sum_{n=-\infty}^{\infty}\delta(\omega-n\Omega)$,
$\Omega=\dfrac{2\pi}{T}$:
\[
    \mathcal{F}[f_T(t)]
    =F_0(j\omega)\,\Omega\,\delta_{\Omega}(\omega)
    =\Omega\sum_{n=-\infty}^{\infty}F_0(jn\Omega)\,\delta(\omega-n\Omega).
\]
Comparing the impulse strengths $2\pi F_n=\Omega T F_n$ with the above,
\[
    F_n=\frac{1}{T}F_0(jn\Omega),
\]
i.e.\ \textbf{the Fourier coefficient of a periodic signal equals the spectrum
of one single pulse, sampled at $n\Omega$, divided by $T$}.

\section{Frequency-Domain Analysis of LTI Systems}

\subsection{Frequency Response}

The Fourier transform decomposes a signal into a sum of complex exponentials
of infinitely many frequencies, so first study the response of a system to a
complex exponential. From $y(t)=h(t)*f(t)$,
\[
    y(t)=\int_{-\infty}^{\infty}h(\tau)f(t-\tau)\,\dd\tau
    =\int_{-\infty}^{\infty}h(\tau)e^{-j\omega\tau}\,\dd\tau\;e^{j\omega t}
    =H(j\omega)e^{j\omega t}.
\]
So the zero-state response to a complex exponential is a complex exponential
of the same frequency, multiplied by the coefficient $H(j\omega)$. Moreover,
since
\[
    f(t)=\frac{1}{2\pi}\int_{-\infty}^{\infty}F(j\omega)e^{j\omega t}\,\dd\omega,
\]
each frequency component is $\dfrac{F(j\omega)\dd\omega}{2\pi}e^{j\omega t}$,
and by linearity
\[
    y(t)=\int_{-\infty}^{\infty}H(j\omega)\frac{F(j\omega)\,\dd\omega}{2\pi}e^{j\omega t}
    =\frac{1}{2\pi}\int_{-\infty}^{\infty}H(j\omega)F(j\omega)e^{j\omega t}\,\dd\omega .
\]
Hence
\[
    Y(j\omega)=H(j\omega)F(j\omega),
\]
which is precisely the \textbf{time-domain convolution theorem}. $H(j\omega)$
is called the \textbf{frequency response} (system function):
\[
    H(j\omega)=\frac{Y(j\omega)}{F(j\omega)} .
\]

\subsection{Frequency-Domain Analysis}

Convolution of signals in the time domain corresponds to multiplication of
their spectra in the frequency domain.

\begin{note}
\textbf{Example}\quad A system is described by
\[
    y'(t)+2y(t)=f(t),
\]
Find the response to the input $f(t)=e^{-t}\varepsilon(t)$.

\textbf{Solution}\quad Transform the time-domain equation into the frequency
domain:
\[
    j\omega Y(j\omega)+2Y(j\omega)=F(j\omega),
\]
so the frequency response is
\[
    H(j\omega)=\frac{Y(j\omega)}{F(j\omega)}=\frac{1}{j\omega+2}.
\]
Since
\[
    F(j\omega)=\frac{1}{j\omega+1},
\]
we have
\[
    Y(j\omega)=\frac{1}{(j\omega+2)(j\omega+1)}
    =\frac{1}{j\omega+1}-\frac{1}{j\omega+2},
\]
and the inverse transform gives
\[
    y(t)=\bigl(e^{-t}-e^{-2t}\bigr)\varepsilon(t).
\]
\end{note}

\subsection{Distortionless Transmission}

Distortionless transmission means the output differs from the input only in
amplitude and time of occurrence, without any change of waveform. With input
$f(t)$,
\[
    y(t)=Kf(t-t_0),
    \qquad
    Y(j\omega)=Ke^{-j\omega t_0}F(j\omega),
\]
so the frequency response must be
\[
    H(j\omega)=Ke^{-j\omega t_0},
\]
with impulse response
\[
    h(t)=K\delta(t-t_0).
\]
That is: \textbf{a constant magnitude response together with a phase response
that is a straight line through the origin} is necessary and sufficient for
distortionless transmission.

\subsection{Response of the Ideal Low-Pass Filter}

An ideal low-pass filter passes signals with angular frequency below
$\omega_c$ without distortion and blocks those above $\omega_c$. The passing
range is the passband and the blocked range the stopband.

Let the cutoff angular frequency be $\omega_c$, with
$|H(j\omega)|=1$ in the passband and phase $\varphi(\omega)=-\omega t_0$;
then
\[
    H(j\omega)=
    \begin{cases}
        e^{-j\omega t_0}, & |\omega|<\omega_c,\\
        0, & |\omega|>\omega_c,
    \end{cases}
\]
which is a gate function of width $2\omega_c$:
\[
    H(j\omega)=e^{-j\omega t_0}g_{2\omega_c}(\omega).
\]

\begin{figure}[htbp]
\centering
\begin{tikzpicture}
\begin{axis}[
    width=7.2cm, height=4.2cm,
    xmin=-4.2, xmax=4.2, ymin=0, ymax=1.35,
    xlabel={$\omega$}, ylabel={$|H(j\omega)|$},
    axis lines=middle,
    xtick={-3,0,3}, xticklabels={$-\omega_c$,$0$,$\omega_c$},
    ytick={0,1},
    title={Magnitude response},
    title style={font=\small},
]
    \addplot[blue, very thick, const plot]
        coordinates {(-4.2,0) (-3,0) (-3,1) (3,1) (3,0) (4.2,0)};
\end{axis}
\end{tikzpicture}
\hspace{6mm}
\begin{tikzpicture}
\begin{axis}[
    width=7.2cm, height=4.2cm,
    xmin=-4.2, xmax=4.2, ymin=-0.9, ymax=0.9,
    xlabel={$\omega$}, ylabel={$\varphi(\omega)$},
    axis lines=middle,
    xtick={-3,0,3}, xticklabels={$-\omega_c$,$0$,$\omega_c$},
    ytick={-0.6,0,0.6},
    title={Phase response (in the passband)},
    title style={font=\small},
]
    \addplot[red, very thick, domain=-3:3] {-0.25*x};
\end{axis}
\end{tikzpicture}
\caption{Frequency characteristics of the ideal low-pass filter}
\end{figure}

\subsubsection{Impulse Response}

\[
    h(t)=\mathcal{F}^{-1}\bigl[e^{-j\omega t_0}g_{2\omega_c}(\omega)\bigr].
\]
By duality: $g_\tau(t)\lrarr\tau\Sa\!\left(\dfrac{\tau}{2}\omega\right)$, so
$\tau\Sa\!\left(\dfrac{\tau}{2}t\right)\lrarr 2\pi g_\tau(-\omega)=2\pi g_\tau(\omega)$.
Taking $\tau=2\omega_c$ and including the time-shift factor,
\[
    h(t)=\frac{\omega_c}{\pi}\,\Sa\bigl[\omega_c(t-t_0)\bigr].
\]

\begin{figure}[htbp]
\centering
\begin{tikzpicture}
\begin{axis}[
    width=11cm, height=5cm,
    xmin=-4, xmax=4, ymin=-0.45, ymax=1.25,
    xlabel={$t$}, ylabel={$h(t)$},
    axis lines=middle,
    xtick={-3,-2,-1,0,1,2,3},
    ytick={0,0.5,1},
    samples=200,
    legend pos=north east,
]
    \addplot[blue, thick, domain=-4:-0.02] {sin(deg(3*x))/(3*x)};
    \addplot[blue, thick, domain=0.02:4]  {sin(deg(3*x))/(3*x)};
    \addlegendentry{$h(t)=\dfrac{\omega_c}{\pi}\Sa(\omega_c t)$}
\end{axis}
\end{tikzpicture}
\caption{Impulse response of the ideal low-pass filter
($\omega_c=3$, $t_0=0$)}
\end{figure}

\subsubsection{Step Response}

Let $g(t)$ be the step response of the ideal low-pass filter; it is the
convolution of $h(t)$ with the unit step:
\[
    g(t)=h(t)*\varepsilon(t)=\int_{-\infty}^{t}h(\tau)\,\dd\tau
    =\int_{-\infty}^{t}\frac{\omega_c}{\pi}\Sa\bigl[\omega_c(\tau-t_0)\bigr]\dd\tau,
\]
\[
    g(t)=\int_{-\infty}^{t}\frac{\omega_c}{\pi}
    \frac{\sin\bigl[\omega_c(\tau-t_0)\bigr]}{\omega_c(\tau-t_0)}\,\dd\tau
    =\frac{1}{2}+\frac{1}{\pi}\Si\bigl[\omega_c(t-t_0)\bigr],
\]
where
\[
    \Si(x)=\int_0^{x}\frac{\sin\eta}{\eta}\,\dd\eta .
\]
The step response starts to build up before $t=t_0$ and still oscillates
afterwards, which shows that the ideal low-pass filter is \textbf{non-causal}
and physically unrealizable; it can only be approximated.

\section{Chapter Summary}

\begin{enumerate}
    \item Everything grows out of \textbf{orthogonality} --- a vanishing inner
          product lets each coefficient be extracted independently.
    \item Series $\to$ transform: as $T\to\infty$, $F_n$ tends to zero and the
          gaps between spectral lines close, so \textbf{density} replaces
          amplitude and \textbf{integration} replaces summation.
    \item The spectrum of a real signal is necessarily \textbf{conjugate
          symmetric}, which is why looking only at the half $\omega\ge0$
          suffices.
    \item Time-domain convolution equals frequency-domain multiplication, so
          LTI system analysis in the frequency domain reduces to the algebraic
          relation $Y=HF$ --- this is the whole value of frequency-domain
          analysis.
\end{enumerate}

\end{document}
